% Generated by tools/edn_latex.py from reports/edn.md.
% Do not edit: change reports/edn.md and run `make edn`.
\ednentry{
  id = {EDN-44},
  title = {The traceability gate is milestone-scoped, and its expected-coverage set is a validated YAML file},
  date = {2026-09-30},
  milestone = {M3: Quality Assurance},
  topic = {Testing Strategy, CI/CD},
  participants = {Lukas},
  decision = {NFR-07's requirement-to-test matrix is gated per milestone rather than over the whole table, and which requirement owes evidence when lives in \texttt{tools/\allowbreak{}expected\_\allowbreak{}coverage.\allowbreak{}yaml}, a schema-validated file the generator checks, rather than in a \texttt{Due} column of the specification's tables. Milestone names are validated against M1 to M6 and the enforced set is read off that known order, never off the order the file lists them in.},
  rationale = {\emph{Alternatives considered:}
\begin{itemize}[leftmargin=*, topsep=2pt]
\item \textbf{Option A (chosen): a small validated YAML file, gate enforced per milestone up to \texttt{current}.}
\newline Pros: the due date of a requirement is a parameter of a CI gate, not a statement about the system, so it does not belong in a requirement document; the generator can enforce a strict schema (every documented requirement booked exactly once, no unknown ID, no unknown or out-of-order milestone) where a table cell cannot be checked at all; bumping \texttt{current} is a one-line, dated, reviewable diff; the file carries the reasoning per milestone in comments.
\newline Cons: a second place to look; the gate is only as honest as the bookings, so a requirement can be deferred by moving one line, which is why the M3 membership and the verification route of every already-due requirement are pinned by a test.
\item \textbf{Option B: a \texttt{Due} column in the specification's tables.}
\newline Pros: everything in one document, visible next to the requirement it belongs to.
\newline Cons: mixes a CI parameter into a document that is supposed to describe the system; \texttt{current} would still have needed a home in prose; the specification's tables are already the widest thing in the docs; nothing can validate a table cell, so a typo or a missing cell would be invisible.
\item \textbf{Option C: no completeness gate at all, matrix published for a human to read.}
\newline Pros: never a red build nobody can fix; zero maintenance.
\newline Cons: leaves NFR-07's promise to a review that nothing reminds anyone to do, which is the state that let the matrix go unimplemented until this PR.
\item \textbf{Option D: enforce completeness over the whole table from the start.}
\newline Pros: no bookkeeping, no bookings to argue about.
\newline Cons: before M4 there is no API, so fourteen of the sixteen functional requirements cannot have a test; the build would be red from the first commit until M6, and a gate that is always red stops being read.
\end{itemize}
\par
\emph{Why:} A gate is only useful if someone acts on it, which rules out D, and only trustworthy if something enforces it, which rules out C. Between A and B, the deciding argument is that the generator can refuse a malformed expected-coverage set and cannot refuse a malformed table cell: the file makes the gate tighten by itself, because a requirement added to the documents fails the build until someone decides when its evidence is due. The decision rests on two specification amendments, both recorded in this PR: NFR-07 said CI ``does not enforce completeness'' and now describes this gate, and the ``Verified by'' legend claimed \textbf{[automated]} means a test already carries the marker, which was false for fourteen of sixteen functional entries and is now described as the route to the evidence with the due date pointing at \texttt{tools/\allowbreak{}expected\_\allowbreak{}coverage.\allowbreak{}yaml}. Neither amendment changes what the system must do; both make the documents describe what CI actually does. An adversarial review of the first implementation found three ways to drop a requirement out of the gate with a one-line edit, one of which -- appending an out-of-order milestone block, because the enforced set was positional in the file -- was not a reviewable statement at all; validating the names against M1 to M6 and reading the order off that closes it, and the other two are now pinned by a test so the argument for loosening the gate has to be written down.},
  ai = {Alternative generation, Alternative assessment, Recommendation, Solution generation},
  response = {Accepted with modifications},
  assessment = {AI laid out A to D with the trade-offs and recommended A, and implemented it. An adversarial AI review of that implementation then found, by execution rather than by reading, that the gate could be loosened three ways without a reviewable statement, and that the enforced set was positional in the YAML. That finding is the modification: the validated milestone order and the tests that pin the gate's parameters were not in the first design and are the reason the decision is defensible as written.},
  aievidence = {Claude Code sessions on 2026-09-30: the options were laid out and A implemented for issue \#40; a second, adversarial review session reproduced each loosening route by running the generator against a modified file and reported the exit codes, and the fixes were then verified the same way.},
  evidence = {\href{https://github.com/mlops-2627q1-mds-upc/MLOPS_Recommenditos/issues/40}{issue \#40}; \href{https://github.com/mlops-2627q1-mds-upc/MLOPS_Recommenditos/pull/51}{PR \#51}; \texttt{tools/\allowbreak{}expected\_\allowbreak{}coverage.\allowbreak{}yaml}; \texttt{DUE\_\allowbreak{}AT\_\allowbreak{}M3} in \texttt{tests/\allowbreak{}test\_\allowbreak{}requirement\_\allowbreak{}matrix.\allowbreak{}py}.},
}
